\documentclass[11pt]{amsart}
\usepackage{amsmath,amssymb,amsthm}
\usepackage[T1]{fontenc}
\usepackage[utf8]{inputenc}
\usepackage{microtype}
\usepackage{hyperref}
\usepackage{enumitem}
\usepackage{booktabs}
\usepackage{longtable}

\hypersetup{
  colorlinks=true,
  linkcolor=blue,
  citecolor=blue,
  urlcolor=blue,
}

\theoremstyle{plain}
\newtheorem{theorem}{Theorem}
\newtheorem{proposition}[theorem]{Proposition}
\newtheorem{corollary}[theorem]{Corollary}

\theoremstyle{definition}
\newtheorem{definition}[theorem]{Definition}
\newtheorem{example}[theorem]{Example}
\newtheorem{remark}[theorem]{Remark}

\title{The element-order collision probability distinguishes the $47$ groups of order $120$}
\author{[Author Name]}
\address{[Affiliation]}
\email{[email@example.org]}
\date{\today}

\subjclass[2020]{Primary 20D60; Secondary 20E45, 05A18}
\keywords{Finite groups, element orders, order type, group invariants, groups of order 120}

\begin{document}

\begin{abstract}
For a finite group $G$, let $N_k(G)$ denote the number of elements of order $k$ in $G$, and define the \emph{element-order collision probability}
\[
  P_2(G) \;=\; \sum_{k \geq 1} \left(\frac{N_k(G)}{|G|}\right)^2.
\]
This is the R\'{e}nyi-$2$ collision entropy of the order-type distribution viewed as a probability mass function on $\mathbb{Z}_{>0}$.
We show by direct enumeration over the GAP \texttt{SmallGroups} library that the $47$ pairwise non-isomorphic groups of order $120$ have $47$ pairwise distinct values of $P_2$, so $P_2$ is a complete invariant on this order. Among these $47$ groups, $27$ ordered pairs $(G_1, G_2)$ satisfy the \emph{single-bucket-swap} property: their order-type distributions differ in exactly two orders $k_1, k_2$ by the same count $n \geq 1$, equivalently, passing from $G_1$ to $G_2$ redistributes $n$ elements from order $k_1$ to order $k_2$. The pair $(S_5, \mathrm{SL}(2,5))$ realises this pattern with $n=24$, $k_1=2$, $k_2=10$, giving the explicit gap $P_2(S_5) - P_2(\mathrm{SL}(2,5)) = 1/300$. All claims are reproducible from a self-contained GAP script of fewer than 50 lines.
\end{abstract}

\maketitle

\section{Introduction}

The multiset of element orders of a finite group $G$, often called the \emph{order type} of $G$ \cite{Thompson1987}, has been intensively studied as a coarse but information-rich invariant. Thompson asked whether sameness of order type forces sameness of solvability; the question was answered negatively by Piwek~\cite{Piwek2024} in 2024. Cameron and Dey~\cite{CameronDey2023} place the order type in a partial order and study its combinatorial structure. The sum of element orders $\psi(G) = \sum_{x \in G} |x|$ is a scalar derived from the order-type distribution and has accumulated a substantial literature since Amiri, Jafarian Amiri and Isaacs~\cite{AmiriAmiriIsaacs2009}; see also~\cite{HerzogLongobardiMaj2018} and references therein.

Less attention has been paid to the second moment of the order-type distribution, equivalently the \emph{collision probability}
\[
  P_2(G) \;=\; \Pr_{(x,y) \in G^2}[|x| = |y|] \;=\; \sum_{k \geq 1} \left(\frac{N_k(G)}{|G|}\right)^2,
\]
where the probability is over two uniformly independent draws from $G$ and $|x|$ denotes the order of $x$. This is the standard collision probability associated with the order-type distribution, equivalently the R\'{e}nyi entropy of order $2$,
\[
  H_2\big(\mathrm{order}(G)\big) \;=\; -\log_2 P_2(G).
\]
The aim of this note is to record a small but striking computational fact: on the order $n = 120$, the scalar $P_2$ is sharp enough to distinguish all $47$ pairwise non-isomorphic groups, and the pairs of groups whose $P_2$ values differ by a particularly small rational gap admit a clean structural description.

\section{Definitions and notation}

Let $G$ be a finite group with $|G| = n$. For $k \in \mathbb{Z}_{>0}$ write
\[
  N_k(G) \;=\; \#\{x \in G : |x| = k\}.
\]
Then $\sum_k N_k(G) = n$ and the \emph{order-type} of $G$ is the multiset $\{|x| : x \in G\}$, equivalently the function $k \mapsto N_k(G)$.

\begin{definition}
The \emph{element-order collision probability} of $G$ is
\[
  P_2(G) \;=\; \sum_{k \geq 1} \left(\frac{N_k(G)}{|G|}\right)^2 \;=\; \frac{1}{|G|^2}\sum_k N_k(G)^2.
\]
\end{definition}

Equivalently $P_2(G)$ is the probability that two uniformly random elements of $G$ have the same order. The denominator $|G|^2$ implies that all $P_2$ values for groups of order $n$ are rationals with denominator dividing $n^2$.

\begin{definition}
Two finite groups $G_1, G_2$ with $|G_1| = |G_2|$ form a \emph{single-bucket-swap pair} with parameters $(k_1, k_2, n)$ if the difference of their order-type distributions is supported on exactly $\{k_1, k_2\}$ with
\[
  N_{k_1}(G_2) - N_{k_1}(G_1) \;=\; -n \quad\text{and}\quad N_{k_2}(G_2) - N_{k_2}(G_1) \;=\; +n,
\]
i.e. passing from $G_1$ to $G_2$ moves $n$ elements from order $k_1$ to order $k_2$ and leaves every other order count unchanged.
\end{definition}

In a single-bucket-swap, the gap in $P_2$ has the closed form
\begin{equation}
  P_2(G_2) - P_2(G_1) \;=\; \frac{1}{|G|^2}\Big[\big(N_{k_2}(G_1) + n\big)^2 - N_{k_2}(G_1)^2 + \big(N_{k_1}(G_1) - n\big)^2 - N_{k_1}(G_1)^2\Big].
\label{eq:swap-gap}
\end{equation}

\section{Main computational results}

\subsection{Distinct values on order 120}

\begin{theorem}\label{thm:distinct-120}
There are $47$ non-isomorphic groups of order $120$. Their $47$ values of $P_2(G)$ are pairwise distinct.
\end{theorem}

\begin{proof}
A direct enumeration over \texttt{AllSmallGroups}$(120)$ in GAP~\cite{GAP4} computes each $N_k(G)$ from the conjugacy class table, then $P_2(G)$ as an exact rational. The $47$ values, sorted ascending, are listed in Table~\ref{tab:p2-values}. They are pairwise distinct. The reproducibility script is given in Section~\ref{sec:reproducibility}.
\end{proof}

\begin{remark}
Theorem~\ref{thm:distinct-120} should not be over-interpreted: $P_2$ is \emph{not} a complete invariant for finite groups in general. Already at order $16$ the order-type alone fails to distinguish, since $\mathbb{Z}_4 \times \mathbb{Z}_4$ and $\mathbb{Z}_2 \times Q_8$ share the same multiset $\{1, 3, 12\}$ (one identity, three involutions, twelve elements of order $4$), hence the same $P_2 = (1 + 9 + 144)/256 = 154/256 = 77/128$. So the smallest order where $P_2$ fails to separate non-isomorphic groups is $\leq 16$, and Theorem~\ref{thm:distinct-120} is a coincidence of order $120$, not a structural fact.
\end{remark}

\subsection{Single-bucket-swap pairs at order 120}

\begin{proposition}\label{prop:swap-count}
Among the $\binom{47}{2} = 1081$ unordered pairs of groups of order $120$, exactly $27$ are single-bucket-swap pairs.
\end{proposition}

\begin{proof}
By direct enumeration; see Table~\ref{tab:swap-pairs} and the GAP script.
\end{proof}

\begin{example}[$S_5$ vs $\mathrm{SL}(2,5)$]
\label{ex:s5-sl25}
The two well-known non-abelian non-solvable groups of order $120$ that contain $A_5$ as a central or quotient factor are $S_5$ and $\mathrm{SL}(2,5)$ (the binary icosahedral group, double cover of $A_5$). Their order-type distributions, computed from their conjugacy class tables, are:
\[
\begin{aligned}
S_5:        \quad & N_1 = 1,\; N_2 = 25,\; N_3 = 20,\; N_4 = 30,\; N_5 = 24,\; N_6 = 20, \\
\mathrm{SL}(2,5): \quad & N_1 = 1,\; N_2 = 1,\; N_3 = 20,\; N_4 = 30,\; N_5 = 24,\; N_6 = 20,\; N_{10} = 24.
\end{aligned}
\]
The two distributions differ in exactly two orders: order $2$ (where $S_5$ has $25$ elements and $\mathrm{SL}(2,5)$ has only the unique central involution) and order $10$ (which has no representatives in $S_5$ but $24$ in $\mathrm{SL}(2,5)$). Passing from $\mathrm{SL}(2,5)$ to $S_5$ redistributes $24$ elements from order $10$ to order $2$. By~\eqref{eq:swap-gap},
\[
  P_2(S_5) - P_2(\mathrm{SL}(2,5)) \;=\; \frac{1}{120^2}\Big[25^2 - 1^2 + 0^2 - 24^2\Big] \;=\; \frac{624 - 576}{14400} \;=\; \frac{48}{14400} \;=\; \frac{1}{300}.
\]
\end{example}

\subsection{The full $P_2$ value table}

\begin{longtable}{rlr@{$\,/\,$}l}
\caption{All $47$ groups of order $120$ sorted by $P_2(G)$ ascending. ``$k$'' is the index in the GAP \texttt{SmallGroups} library at order $120$.}\label{tab:p2-values}\\
\toprule
$k$ & Structure & \multicolumn{2}{c}{$P_2$} \\
\midrule
\endfirsthead
\toprule
$k$ & Structure & \multicolumn{2}{c}{$P_2$} \\
\midrule
\endhead
\midrule
\multicolumn{4}{r}{\textit{(continued on next page)}}\\
\endfoot
\bottomrule
\endlastfoot
17 & $C_{12} \times D_{10}$ & 181 & 1440 \\
4 & $C_{120}$ & 187 & 1440 \\
20 & $C_{3} \times ((C_{10} \times C_{2}) \rtimes C_{2})$ & 43 & 288 \\
31 & $C_{60} \times C_{2}$ & 221 & 1440 \\
22 & $C_{20} \times S_{3}$ & 391 & 2400 \\
10 & $C_{3} \rtimes (C_{4} \times D_{10})$ & 1181 & 7200 \\
27 & $C_{4} \times D_{30}$ & 83 & 480 \\
2 & $C_{3} \times (C_{5} \rtimes C_{8})$ & 251 & 1440 \\
32 & $C_{15} \times D_{8}$ & 17 & 96 \\
11 & $C_{3} \rtimes ((C_{10} \times C_{2}) \rtimes C_{2})$ & 1303 & 7200 \\
8 & $(C_{3} \rtimes C_{4}) \times D_{10}$ & 1321 & 7200 \\
6 & $C_{3} \times (C_{5} \rtimes C_{8})$ & 89 & 480 \\
25 & $C_{5} \times ((C_{6} \times C_{2}) \rtimes C_{2})$ & 1343 & 7200 \\
18 & $C_{3} \times D_{40}$ & 271 & 1440 \\
40 & $C_{6} \times (C_{5} \rtimes C_{4})$ & 277 & 1440 \\
13 & $C_{3} \rtimes ((C_{10} \times C_{2}) \rtimes C_{2})$ & 1403 & 7200 \\
15 & $C_{5} \times SL(2,3)$ & 1411 & 7200 \\
30 & $C_{3} \rtimes ((C_{10} \times C_{2}) \rtimes C_{2})$ & 283 & 1440 \\
19 & $C_{6} \times (C_{5} \rtimes C_{4})$ & 19 & 96 \\
5 & $\mathrm{SL}(2,5)$ & 1427 & 7200 \\
36 & $S_{3} \times (C_{5} \rtimes C_{4})$ & 1429 & 7200 \\
34 & $S_{5}$ & 1451 & 7200 \\
16 & $C_{3} \times (C_{5} \rtimes Q_{8})$ & 97 & 480 \\
35 & $C_{2} \times A_{5}$ & 1457 & 7200 \\
12 & $C_{3} \rtimes D_{40}$ & 1471 & 7200 \\
1 & $C_{5} \times (C_{3} \rtimes C_{8})$ & 493 & 2400 \\
9 & $S_{3} \times (C_{5} \rtimes C_{4})$ & 1493 & 7200 \\
43 & $C_{10} \times A_{4}$ & 1513 & 7200 \\
37 & $C_{5} \times S_{4}$ & 1547 & 7200 \\
33 & $C_{15} \times Q_{8}$ & 323 & 1440 \\
38 & $C_{5} \rtimes S_{4}$ & 1619 & 7200 \\
14 & $C_{3} \rtimes (C_{5} \rtimes Q_{8})$ & 1631 & 7200 \\
24 & $C_{10} \times (C_{3} \rtimes C_{4})$ & 1649 & 7200 \\
23 & $C_{5} \times D_{24}$ & 187 & 800 \\
39 & $A_{4} \times D_{10}$ & 563 & 2400 \\
44 & $C_{2} \times C_{6} \times D_{10}$ & 23 & 96 \\
42 & $C_{2} \times S_{3} \times D_{10}$ & 1813 & 7200 \\
21 & $C_{5} \times (C_{3} \rtimes Q_{8})$ & 629 & 2400 \\
3 & $C_{3} \rtimes (C_{5} \rtimes C_{8})$ & 137 & 480 \\
28 & $D_{120}$ & 47 & 160 \\
47 & $C_{30} \times C_{2} \times C_{2}$ & 85 & 288 \\
7 & $C_{3} \rtimes (C_{5} \rtimes C_{8})$ & 427 & 1440 \\
26 & $C_{3} \rtimes (C_{5} \rtimes Q_{8})$ & 29 & 96 \\
41 & $C_{2} \times (C_{15} \rtimes C_{4})$ & 437 & 1440 \\
29 & $C_{2} \times (C_{15} \rtimes C_{4})$ & 89 & 288 \\
45 & $C_{2} \times C_{10} \times S_{3}$ & 2261 & 7200 \\
46 & $C_{2} \times C_{2} \times D_{30}$ & 481 & 1440 \\

\end{longtable}

\begin{longtable}{rrrrrr@{$\,/\,$}l}
\caption{All $27$ single-bucket-swap pairs at order $120$, sorted first by $n$ (number of elements moved) and then by denominator of the $P_2$ gap. Each row records that passing from group $A$ to group $B$ moves $n$ elements from order $k_1$ to order $k_2$.}\label{tab:swap-pairs}\\
\toprule
$A$ & $B$ & $n$ & $k_1$ & $k_2$ & \multicolumn{2}{c}{$|P_2(A) - P_2(B)|$} \\
\midrule
\endfirsthead
\toprule
$A$ & $B$ & $n$ & $k_1$ & $k_2$ & \multicolumn{2}{c}{$|P_2(A) - P_2(B)|$} \\
\midrule
\endhead
\midrule
\multicolumn{7}{r}{\textit{(continued on next page)}}\\
\endfoot
\bottomrule
\endlastfoot
8 & 41 & 24 & 20 & 4 & 3 & 25 \\
5 & 34 & 24 & 10 & 2 & 1 & 300 \\
10 & 14 & 30 & 2 & 4 & 1 & 16 \\
8 & 12 & 30 & 4 & 2 & 1 & 48 \\
9 & 13 & 30 & 4 & 2 & 1 & 80 \\
29 & 30 & 30 & 4 & 2 & 9 & 80 \\
30 & 46 & 30 & 4 & 2 & 11 & 80 \\
5 & 35 & 30 & 4 & 2 & 1 & 240 \\
11 & 42 & 30 & 4 & 2 & 17 & 240 \\
26 & 27 & 30 & 4 & 2 & 31 & 240 \\
27 & 28 & 30 & 4 & 2 & 29 & 240 \\
2 & 3 & 40 & 24 & 8 & 1 & 9 \\
6 & 7 & 40 & 24 & 8 & 1 & 9 \\
19 & 29 & 40 & 12 & 4 & 1 & 9 \\
40 & 41 & 40 & 12 & 4 & 1 & 9 \\
16 & 26 & 40 & 12 & 4 & 1 & 10 \\
18 & 28 & 40 & 6 & 2 & 19 & 180 \\
44 & 46 & 40 & 6 & 2 & 17 & 180 \\
1 & 3 & 48 & 40 & 8 & 2 & 25 \\
21 & 26 & 48 & 20 & 4 & 1 & 25 \\
24 & 29 & 48 & 20 & 4 & 2 & 25 \\
23 & 28 & 48 & 10 & 2 & 3 & 50 \\
45 & 46 & 48 & 10 & 2 & 1 & 50 \\
29 & 46 & 60 & 4 & 2 & 1 & 40 \\
3 & 26 & 60 & 8 & 4 & 1 & 60 \\
3 & 28 & 60 & 8 & 2 & 1 & 120 \\
26 & 28 & 60 & 4 & 2 & 1 & 120 \\

\end{longtable}

\section{Discussion}

The result of Theorem~\ref{thm:distinct-120} is a discrete coincidence: at order $120$, the $47$-element image of $P_2$ in $\mathbb{Q}$ happens to be injective. There is no reason to expect this on every order. Indeed, the order-type itself fails to be a complete invariant from order $16$ onward, and any function of the order-type inherits this limitation in the worst case.

A natural question is to characterise the orders $n$ at which $P_2$ is injective on the set of isomorphism classes. Computer experiments accessible to a workstation can answer this for $n \leq 256$ or so; we leave the systematic computation, and the question of structural reasons (if any) for injectivity at specific orders, to follow-up work.

The single-bucket-swap pattern of Proposition~\ref{prop:swap-count} is a rigidity condition: most pairs of groups of order $120$ differ in three or more order-buckets. The $27$ that differ in exactly two are a small minority and tend to occur between groups that are structurally close (sharing a normal subgroup, or being central extensions of the same quotient by isomorphic kernels). Whether ``single-bucket-swap'' admits an algebraic characterisation in terms of standard group-theoretic relations is open.

\section{Reproducibility}
\label{sec:reproducibility}

All numerical claims are reproducible from the GAP script \texttt{reproducibility.g}, which we list in full below. On a workstation the entire computation runs in under one second.

\begin{verbatim}
LoadPackage("smallgrp");

n := 120;
num := NumberSmallGroups(n);
results := [];

for i in [1..num] do
    G := SmallGroup(n, i);
    desc := StructureDescription(G);
    ccs := ConjugacyClasses(G);
    orders := List(ccs, c -> Order(Representative(c)));
    sizes  := List(ccs, Size);
    distinct_orders := Set(orders);
    Nk := [];
    for k in distinct_orders do
        nk := Sum(Filtered([1..Length(orders)], j -> orders[j] = k),
                  j -> sizes[j]);
        Add(Nk, [k, nk]);
    od;
    p2 := Sum(Nk, x -> (x[2]/n)^2);
    Add(results, rec(id := i, desc := desc, Nk := Nk, p2 := p2));
od;

# Theorem: distinct P_2
distinct_p2 := Set(List(results, r -> r.p2));
Print("Distinct P_2 count: ", Length(distinct_p2), "\n");

# Proposition: 27 single-bucket-swap pairs
swap := 0;
for i in [1..num] do for j in [(i+1)..num] do
    a := results[i].Nk;; b := results[j].Nk;;
    all_k := Union(List(a, x -> [x[1]]), List(b, x -> [x[1]]));
    diff := [];
    for k in all_k do
        na := First(a, x -> x[1] = k); if na = fail then na := [k, 0]; fi;
        nb := First(b, x -> x[1] = k); if nb = fail then nb := [k, 0]; fi;
        if na[2] <> nb[2] then Add(diff, [k, nb[2] - na[2]]); fi;
    od;
    if Length(diff) = 2 and diff[1][2] + diff[2][2] = 0 then
        swap := swap + 1;
    fi;
od; od;
Print("Single-bucket-swap pair count: ", swap, "\n");
\end{verbatim}

A worked Python post-processing script (parsing the GAP output, sorting the $P_2$ values, and listing the $27$ swap pairs) is included in the supplementary materials accompanying this preprint.

\begin{thebibliography}{99}

\bibitem{AmiriAmiriIsaacs2009}
H. Amiri, S.~M. Jafarian Amiri, and I.~M. Isaacs,
\emph{Sums of element orders in finite groups},
Comm. Algebra \textbf{37} (2009), no. 9, 2978--2980.

\bibitem{CameronDey2023}
P.~J. Cameron and H.~K. Dey,
\emph{On the order sequence of a group},
arXiv:2310.06516 (2023), to appear in Electron. J. Combin.

\bibitem{GAP4}
The GAP Group,
\emph{GAP -- Groups, Algorithms, and Programming},
Version 4.13, 2024. \url{https://www.gap-system.org}.

\bibitem{HerzogLongobardiMaj2018}
M. Herzog, P. Longobardi, and M. Maj,
\emph{An exact upper bound for sums of element orders in non-cyclic finite groups},
J. Pure Appl. Algebra \textbf{222} (2018), no. 7, 1628--1642.

\bibitem{Piwek2024}
J. Piwek,
\emph{Solvable and non-solvable finite groups of the same order type},
arXiv:2403.02197 (2024), Algebra \& Number Theory \textbf{19} (2025).

\bibitem{Thompson1987}
J.~G. Thompson,
\emph{Unpublished letter, $1987$};
see Piwek~\cite{Piwek2024} for context.

\end{thebibliography}

\end{document}
